\chapter{From an English Instruction to Shuttle Motion}
\label{chap:neurosymbolic}

\section{Runtime architecture}

\Cref{fig:neurosymbolic-architecture} follows English and camera inputs through
validated proposals, PlanSys2/POPF planning, supervised typed commands and
verified Gazebo effects; an incomplete goal re-enters observation and planning.
Interface rules are in \cref{chap:interfaces}.

\begin{figure}[p]
  \centering
  \begin{tikzpicture}[
  font=\sffamily\footnotesize,
  execute at begin node={\hyphenpenalty=10000\exhyphenpenalty=10000},
  input/.style={draw=black!42,rounded corners=2pt,fill=black!3,
    align=center,text width=34mm,minimum height=8mm,inner sep=1.2mm},
  learned/.style={draw=mfjablue,very thick,rounded corners=2pt,
    fill=mfjalight,align=center,text width=34mm,minimum height=14mm,
    inner sep=1.3mm},
  visualmodel/.style={learned,text width=72mm,minimum height=16mm,
    inner xsep=1.3mm,inner ysep=1mm},
  accepted/.style={draw=mfjagreen,very thick,rounded corners=2pt,
    fill=mfjagreen!7,align=center,text width=72mm,minimum height=20mm,
    inner sep=1.3mm},
  gateway/.style={draw=mfjagreen,very thick,rounded corners=2pt,
    fill=mfjagreen!7,align=center,text width=154mm,minimum height=19mm,
    inner sep=1.4mm},
  loopstage/.style={draw=mfjagreen,very thick,rounded corners=2pt,
    fill=mfjagreen!7,align=center,text width=112mm,minimum height=11mm,
    inner sep=1.2mm},
  planner/.style={draw=mfjaorange,very thick,rounded corners=2pt,
    fill=mfjaorange!8,align=center,text width=112mm,minimum height=12mm,
    inner sep=1.2mm},
  supervisor/.style={draw=mfjared,very thick,rounded corners=2pt,
    fill=mfjared!6,align=center,text width=112mm,minimum height=12mm,
    inner sep=1.2mm},
  controller/.style={draw=black!58,very thick,rounded corners=2pt,
    fill=black!5,align=center,text width=112mm,minimum height=14mm,
    inner sep=1.2mm},
  verify/.style={draw=mfjagreen,very thick,rounded corners=2pt,
    fill=mfjagreen!7,align=center,text width=112mm,minimum height=14mm,
    inner sep=1.3mm},
  success/.style={draw=mfjagreen,very thick,rounded corners=8pt,
    fill=mfjagreen!13,align=center,text width=61mm,minimum height=9mm,
    inner sep=1mm,font=\sffamily\footnotesize\bfseries},
  replan/.style={draw=mfjagreen,very thick,densely dashed,rounded corners=8pt,
    fill=white,align=center,text width=61mm,minimum height=9mm,
    inner sep=1mm,font=\sffamily\footnotesize\bfseries,
    text=mfjagreen!75!black},
  paneltitle/.style={anchor=west,font=\sffamily\small\bfseries,
    text=mfjablue},
  flow/.style={-{Latex[length=2.1mm]},very thick,draw=mfjablue},
  feedback/.style={-{Latex[length=2.1mm]},very thick,densely dashed,
    draw=mfjagreen}
]
  % Panel A: the confirmed task and the visual path meet in
  % the task gateway.  Every label is at least footnotesize for A4 printing.
  \node[paneltitle] at (-79mm,10mm)
    {A \enspace INTERPRET THE REQUEST AND ACCEPT THE SCENE STATE};
  \draw[mfjablue!28,line width=0.55pt] (-79mm,7mm) -- (79mm,7mm);

  \node[input,text width=40mm] (english) at (-43mm,0) {
    \textbf{English instruction}\\task-goal CLI};
  \node[learned] (qwen) at (-62mm,-19mm) {
    \textbf{Pretrained Qwen2.5-1.5B}\\
    Instruct; local strict-JSON proposal};
  \node[learned,draw=mfjagreen,fill=mfjagreen!7] (facts)
    at (-24mm,-19mm) {
    \textbf{Explicit-fact extractor}\\
    explicit fields from the same request};
  \node[accepted] (task) at (-43mm,-42mm) {
    \textbf{Combine model and explicit fields}\\[-0.2ex]
    field fusion $\rightarrow$ \texttt{TaskGoalDraft}\\
    validate, clarify if needed and confirm $\rightarrow$ \texttt{TaskGoal}};

  \node[input,text width=72mm] (visualinputs) at (42mm,0) {
    \textbf{Paired camera inputs}\\synchronised left/right RGB frames +
    shuttle identity, presence and timestamps};
  \node[visualmodel] (vfour) at (42mm,-21mm) {
    \textbf{Split-rail visual model}\\[-0.2ex]
    shared low-level stem + rail-isolated left/right branches};
  \node[accepted] (visual) at (42mm,-45mm) {
    \textbf{Apply frozen scores and validate the complete frame}\\[-0.2ex]
    segment/loaded score floors + freshness/contract checks\\
    $\rightarrow$ accepted \texttt{VisualStateObservation} (declared schema)};

  \node[gateway] (ground) at (0,-70mm) {
    \textbf{Task gateway: ground the goal and build a fresh PDDL problem}\\[-0.2ex]
    confirmed \texttt{TaskGoal} + accepted visual observation + presence/device/DZI state
    $\rightarrow$ validated \texttt{ObservedState} and current symbolic facts\\[-0.2ex]
    \textit{Optional direct ProblemExpert predicate mirror: disabled}};

  \draw[flow] ([xshift=-19mm]english.south) -- (qwen.north);
  \draw[flow] ([xshift=19mm]english.south) -- (facts.north);
  \draw[flow] (qwen.south) -- ([xshift=-19mm]task.north);
  \draw[flow] (facts.south) -- ([xshift=19mm]task.north);
  \draw[flow] (visualinputs) -- (vfour);
  \draw[flow] (vfour) -- (visual);
  \draw[flow] (task.south) -- ([xshift=-43mm]ground.north);
  % Keep the final segment vertical so the arrowhead points unambiguously from
  % the accepted visual frame into the task gateway.
  \draw[flow] (visual.south) -- ([xshift=42mm]ground.north);

  % Panel B: one-step receding-horizon planning and result checking.
  \node[paneltitle] at (-79mm,-85mm)
    {B \enspace PLAN, EXECUTE ONE SAFE STEP, AND CHECK THE RESULT};
  \draw[mfjablue!28,line width=0.55pt] (-79mm,-88mm) -- (79mm,-88mm);

  \node[planner] (plan) at (0,-99mm) {
    \textbf{PlanSys2 planner service / POPF}\\[-0.2ex]
    current domain + gateway-built problem $\rightarrow$ symbolic plan};
  \node[loopstage] (translate) at (0,-115mm) {
    \textbf{Translate and select the first supported step}\\[-0.2ex]
    first plan step $\rightarrow$ sparse named live command};
  \node[supervisor] (supervise) at (0,-131mm) {
    \textbf{Deterministic safety supervisor}\\[-0.2ex]
    verify the sparse command $\rightarrow$ publish a typed ROS~2 command};
  \node[controller] (controller) at (0,-148mm) {
    \textbf{Kinematic rail controller and Gazebo}\\[-0.2ex]
    apply the command $\rightarrow$ shuttle motion; return controller state
    and identity-bearing DZI readings};
  \node[verify] (verify) at (0,-167mm) {
    \textbf{Wait for newer evidence and compare the observed effect}\\[-0.2ex]
    accepted visual observation + supervisor/controller state + DZI evidence};
  \node[success] (success) at (-35mm,-186mm)
    {complete $\rightarrow$ \texttt{task\_goal\_satisfied}};
  \node[replan] (replan) at (35mm,-186mm)
    {not complete $\rightarrow$ re-observe and replan};

  % Enter Panel B at the right margin so the arrow does not cross its title;
  % finish vertically into the planner so the direction remains unambiguous.
  \draw[flow] ([xshift=72mm]ground.south) -- (72mm,-90mm)
    -- (50mm,-90mm) -- ([xshift=50mm]plan.north);
  \draw[flow] (plan) -- (translate);
  \draw[flow] (translate) -- (supervise);
  \draw[flow] (supervise) -- (controller);
  \draw[flow] (controller) -- (verify);
  \draw[flow,draw=mfjagreen] (verify.south) -- (success.north);
  \draw[feedback] (verify.south) -- (replan.north);

  % Only an incomplete task returns to a post-command accepted observation.
  % Keep a dedicated right-hand lane and a final segment longer than the arrow
  % head; the previous near-zero segment made the arrowhead fold at the corner.
  \coordinate (replanLaneBottom) at (81.2mm,-186mm);
  \coordinate (replanLaneTop) at (81.2mm,-45mm);
  \draw[
    feedback,
    -{Latex[length=1.4mm,width=1.15mm]},
    line join=round
  ] (replan.east) -- (replanLaneBottom)
    -- (replanLaneTop) -- (visual.east);

  % Explicit extent avoids global scaling, preserving the footnotesize text.
  \path[use as bounding box] (-82mm,14mm) rectangle (82mm,-196mm);
\end{tikzpicture}

  \caption[Room~315 English-to-motion runtime architecture.]
  {End-to-end sequence from task and paired-camera inputs to a verified Gazebo
  result; the dashed green path denotes observation and replanning while the
  grounded goal remains incomplete.}
  \label{fig:neurosymbolic-architecture}
\end{figure}
\FloatBarrier

\section{Step 1---interpret and confirm the English task}

The CLI applies \cref{chap:ai-development}, clarifies missing fields and requires
confirmation before publishing \code{TaskGoal}. It resolves explicit identities
and colour aliases; scene-dependent policies remain ungrounded. \emph{Grounding}
selects one shuttle and, for station requests, an admissible slot; the planner,
not the language model, chooses actions.

\section{Step 2---read the current scene and ground the goal}

The visual runtime verifies its manifest---the hash-bound artefact, schema,
configuration and Gazebo-scope record---plus checkpoint, preprocessing and
camera order. Any missing, stale, conflicting or low-confidence requirement
rejects the pair; an accepted \code{VisualStateObservation} carries schema and
provenance. The inference node does not mutate ProblemExpert: the gateway fuses
the observation with presence, device and identity-bearing sensor reports into
\code{ObservedState} for planning.

Symbolic shuttle locations are derived directly from the accepted visual
observation rather than from controller coordinates.

The gateway validates and grounds \emph{any}, \emph{nearest} and station-slot
policies against current state/occupancy. Payload-based selection requires, by
default, five distinct consecutive accepted observations within freshness
limits. The grounded identity retains the goal ID and stays fixed for execution.

\section{Step 3---build a fresh PDDL problem and plan}

\subsection{From current state to decision predicates}

Each cycle builds a fresh PDDL problem from the grounded goal and latest
\code{ObservedState}: present shuttles, locations, occupancy, payload and device
facts. PDDL receives no raw image/text; unknown, stale or conflicting state
blocks generation rather than implying free space.

\subsection{How the current facts constrain the next action}

Transport actions require \code{validated\_state}, matching
\code{active\_goal\_side}, and action-specific location, occupancy and route
predicates. POPF minimises dimensionless symbolic cost (setup/relocation
penalties, not travel time) among goal-reaching plans; supervisor checks remain
separate. \Cref{tab:pddl-decision-rules} summarises the principal cases.

In the table, \code{validated\_state} marks contract-valid facts. The field
\code{active\_goal\_side} names the grounded rail. \emph{Normal route} is the
exterior-loop switch arrangement; a \emph{topology block} is a segment-level,
non-slot location; \emph{clearance mode} temporarily relocates blockers;
\emph{capacity-safe} respects holding/separation limits; and a device group is
one route's configured switches/stoppers.

\begin{table}[H]
\centering
\footnotesize
\caption{Principal symbolic decision cases in the current Room~315 PDDL domain.}
\label{tab:pddl-decision-rules}
\begin{tabularx}{\textwidth}{@{}
  >{\raggedright\arraybackslash}p{0.22\textwidth}
  >{\raggedright\arraybackslash}p{0.36\textwidth}
  >{\raggedright\arraybackslash}X@{}}
\toprule
\textbf{Current symbolic situation} & \textbf{Decisive conditions} &
\textbf{Applicable action and symbolic result}\\
\midrule
Canonical exterior-route preparation & Validated active side, connected
stations, normal-route and named device-group facts &
\path{prepare_switches}, then \path{open_stoppers}, establish
\code{switches\_ready}, \code{stoppers\_open} and \code{path\_ready} when those
facts are still required\\
\tablerowrule
Direct move from an exact slot & Source occupancy agrees with the selected
shuttle; target is \code{slot\_free}; route is clear and ready; no clearance is
pending & \path{move_shuttle_to_slot} releases the source occupancy and marks
the target occupied and reserved\\
\tablerowrule
Switch-specific or segment-origin route & A topology block is known; a route to
the free target is available and clear; no clearance is pending &
\path{prepare_topology_route} or \path{prepare_slot_topology_route}, followed by
the matching topology move, configures a legal route without inventing a
source slot\\
\tablerowrule
Another shuttle blocks the required route & An ordered clearance is pending;
the next relocation or staging destination is capacity-safe and its interior
entry is clear & The clearance family enters \code{clearance\_mode} and selects
one permitted next relocation. Each completed step produces a fresh problem;
finish or pause actions control route restoration\\
\tablerowrule
Mixed route requiring normalisation & Route reconfiguration is allowed;
clearance mode is inactive; no relocation remains &
\path{restore_normal_route} restores the exterior route and its readiness facts\\
\tablerowrule
Arrival or inspection is ready to close & Target station/slot and \code{DISABLED} controller-state
facts hold, or a confirmed inspection is required & \path{stop_shuttle} and the
terminal actions establish task completion; \path{inspect_state} completes an
inspection without an actuation command\\
\bottomrule
\end{tabularx}
\end{table}

PDDL effects are predictions; the next cycle always uses the observed
post-action state. The executive calls the POPF-backed PlanSys2 planner service
(not its Executor), validates the plan and selects its first supported
\code{PddlPlanStep}/\code{TranslatedPlanStep}: an action the translator can map
to a checked command tied to the planning state.

\section{Step 4---translate, supervise and move}

The translator makes a sparse named \code{TranslatedPlanStep}; the supervisor
applies \cref{chap:interfaces} before publishing a typed device/shuttle command,
after which the kinematic controller updates shuttle state and Gazebo pose.

\section{Step 5---check the result and decide what follows}

The executive checks postconditions against newer state and, if incomplete,
replans from it. Final-slot success requires the DZI/controller conditions in
\cref{chap:interfaces}.

\section{Case B03 from instruction to arrival}

For \emph{Move the green shuttle to slot 4 on the right rail}, the campaign
fixed request, goal and launch state but left planning open. From a cold start
(fresh Gazebo process with no retained task, observation or controller state),
the planner chose one direct step to R3's slot~4/DZI4R destination
(\cref{tab:green-shuttle-trace}).

\begin{table}[H]
\centering
\small
\caption{English-to-motion record for campaign case B03.}
\label{tab:green-shuttle-trace}
\begin{tabularx}{\textwidth}{@{}>{\raggedright\arraybackslash}p{0.21\textwidth}X@{}}
\toprule
\textbf{Stage} & \textbf{Recorded result}\\
\midrule
Interactive request & \emph{Move the green shuttle to slot 4 on the right
rail.} The harness supplied \code{yes} after the rendered interpretation
matched the case declaration\\
\tablerowrule
Confirmed task & Transport
\code{room315\_right\_shuttle\_3} on the right rail to slot~4\\
\tablerowrule
Current state & Visual state supplies current location and
occupancy; a fresh PDDL problem is sent to POPF and the executive selects the
first supported step\\
\tablerowrule
Execution step & The recorded initial state places R3 at slot~3 with
slot~4 free. The motion portion contains one
\code{move\_shuttle\_to\_slot} step; it requires neither blocker clearance nor
a topology relocation\\
\tablerowrule
Motion command & \code{move\_shuttle\_to\_slot} for R3 from right
slot~3 to right slot~4 within the Stäubli TX2-60L section at \SI{0.2}{m/s}\\
\tablerowrule
Arrival & Two separately timestamped DZI4R readings report identity R3;
the matching \code{ShuttleState} reports target slot~4 and \code{DISABLED}, and
a newer visual observation records the final scene\\
\tablerowrule
Outcome & Every recorded postcondition is satisfied; the final status is
\code{succeeded} with reason \code{task\_goal\_satisfied}\\
\bottomrule
\end{tabularx}
\end{table}

Other identity, colour and payload policies use the same workflow, grounding
from the current scene before PDDL construction. Results are in
\cref{chap:results}.
